% =============================================================================
% 07_discussion.tex  --  Chapter 7, Discussion (target 5 pages after the 19 Oct trimming pass). Owner: Discussion writer.
% Plan: docs/thesis/REWRITE_PLAN.md sections 1.8, 2 ("7 Discussion"), 4 (T1), 5 (writing rules).
% Sources: docs/thesis_requirements_plan.md sections 1-3 (status of 18 Oct evening); docs/handover.md 0a, 0b, 5
% (item 7), 6 (corrections log); docs/deck_explained_2026-10-19.md parts 7-8; chapters 3-6 and their tables;
% data/plot_stats.md.
% Numbers: macros of numbers_results / _diagnostics / _method / _data; the few numbers first quoted here are in
% numbers_discussion.tex. Table owned here: tables/T01_requirements.tex (tab:requirements).
% Trimming pass 19 Oct: Table 7.1 carries status and evidence, the requirement paragraphs keep only what is missing;
% the long version is in commit 76aa1a4.
% =============================================================================

\chapter{Discussion}
\label{ch:discussion}

This chapter reads the results against the requirements, explains why the headline error barely moved, states the
limits of each claim and records what the work taught about building such a pipeline.

% -----------------------------------------------------------------------------
\section{The requirements, one by one}
\label{sec:disc-requirements}

The supervisor set three tiers (Section~\ref{sec:intro-requirements}): requirements without which he would not
consider the thesis ready for submission (1a, 1b, 1c), requirements for an acceptable grade (2a, 2b), and a
well-written report with verifiable code for the best grade (3). Table~\ref{tab:requirements} gives the status of
each on \DiscStatusDate{} and its evidence. I split 1a into pose and shape, because their evidence differs. The
paragraphs below name only what is missing and where Chapter~\ref{ch:conclusion} plans it.

% T1 Requirements map (Discussion 7.1, full form; short form in Introduction 1.3).
% Owner: skeleton / Discussion writer. Source: REWRITE_PLAN.md section 4 (T1);
% docs/thesis_requirements_plan.md sections 1-2 (P 1, P 2); deck slide 23 (D slide 23).
% Status words only: met / met with caveats / partly / not met / not started (plan section 5).
% Job ids for the evidence: 3611336:224-229 (3DPW test), 3612131 (RH final), 3612078 (demo),
% 3613616 (RH masked), 3614191 (student, pending), 3614192 (NAVER), 3611284 (label check v23).
% Section numbers in the evidence column are of the NEW thesis; replace by \ref once the
% writers' labels exist.
\begin{table}[tbp]
\centering\footnotesize
\setlength{\tabcolsep}{3pt}
\begin{tabular}{@{}>{\raggedright}p{1.3cm}>{\raggedright}p{3.0cm}>{\raggedright}p{1.5cm}>{\raggedright}p{3.6cm}>{\raggedright\arraybackslash}p{3.3cm}@{}}
\toprule
Req. & Supervisor's wording (short) & Status & Evidence & What is missing \\
\midrule
1a-pose & ``Pose of all body parts must look correct'' on the 3DPW test set and on any in-the-wild image &
  partly &
  3DPW test (v23 labels): PVE \FinalThreeDPW\,mm, recall \FinalRecallThreeDPW\,\%, precision \FinalPrecThreeDPW\,\%
  (Table~\ref{tab:final}); necks and fingers corrected on the holdout (Sec.~\ref{sec:res-bone}, \ref{sec:res-fingers});
  real photos judged by eye only &
  predicted fingers on a real photo; depth slope 0.42; one run, no seed \\
\addlinespace
1a-shape & ``Mesh size, proportions, age, muscles \ldots{} must visibly change according to the person'' &
  partly, at risk &
  Relative Human: gender \RHGenderFinal\,\% (Base \RHGenderBase), child AUC \RHAUCFinal{} (Base \RHAUCBase)
  (Table~\ref{tab:relhuman}); muscle, weight, cup size, firmness near 0.5 on photos; masked arm trades
  variety for gender (Sec.~\ref{sec:res-mask}) &
  per-phenotype correlation on AnnyOne; rest-pose line-up (F17); the partial mask \\
\addlinespace
1b & ``Must not be exactly the same as the Multi-HMR model'': new and better, or significantly smaller and
  $\geq$\,95\,\% as good &
  not met &
  the architecture is Multi-HMR's with an Anny head; Student \StudentParams\,M vs \FinalParams\,M params
  (Table~\ref{tab:baseline}); accuracy pending (Table~\ref{tab:student}); Multi-HMR-Anny (NAVER)
  \NAVERAnnyOne{} / \NAVERThreeDPW\,mm &
  the Student's accuracy; the supervisor's choice of baseline for the bar \\
\addlinespace
1c & Evaluation on the real-world set with Anny parameters, not only SMPL-X &
  met with caveats &
  3DPW test in Anny space, 8,190 images / 11,854 people; label check WARN, body joints 14.1\,mm (Ch.~3) &
  state the label caveats: 8 constant phenotypes, hands and head at rest, $\sim$14\,mm floor \\
\addlinespace
2a & Convert 3DPW train and test to Anny and include the train set in training &
  met &
  train 5,882 images / 7,584 people; 30\,\% of each batch; Control \ControlThreeDPW{} vs Base
  \BaseThreeDPW\,mm on 3DPW test (Table~\ref{tab:final}) &
  nothing (a 50\,\% mix is optional) \\
\addlinespace
2b & The model understands motion in video instead of predicting every frame on its own &
  not started &
  no temporal code &
  design and evaluation (Ch.~\ref{ch:conclusion}) \\
\addlinespace
3 & A well-written report; a clean codebase and models he can verify &
  partly &
  this thesis; Appendix~\ref{app:reproduction} maps every number to its run &
  code release and licences (Ch.~\ref{ch:conclusion}) \\
\bottomrule
\end{tabular}
\caption{The supervisor's requirements and their status. Numbers: AnnyOne holdout in setting V1, 3DPW test on the
corrected (v23) labels, Relative Human validation split; run names as defined in Chapter~4.}
\label{tab:requirements}
\end{table}


\paragraph{1a-pose (partly).} Missing are correct near and far depth (slope \FinalSlope;
Figure~\ref{fig:depth-scatter}), a pose number on real photos, where placement offsets, ghost bodies and hard poses
such as the seated guitar player remain, and a second run. Section~\ref{sec:future} plans a repeat of Final; the slope
stays an open limitation.

\paragraph{1a-shape (partly).} Gender and age change with the person; muscle, weight, cup size and firmness do not,
because Final learned the constants the converter writes for every 3DPW person, and neither mask arm freed them
without losing gender (Section~\ref{sec:res-mask}). Missing are a measure of how well each phenotype follows its label
on AnnyOne, which Section~\ref{sec:future} plans, and a recipe for that variety.

\paragraph{1b (not met).} Final is Multi-HMR with its output resized to Anny, and \DiscFinalSmallerPct\,\% smaller
than NAVER's model, which is not ``significantly smaller''. Missing are the Student's accuracy
\TODO{Student result: update 1b, T1 and the status word}, the supervisor's choice of the reference model, and his
view on whether a Multi-HMR-style network with a smaller backbone counts. If the Student misses while its error still
falls, Section~\ref{sec:future} plans to continue its training.

\paragraph{1c (met with caveats).} The caveats are those of Section~\ref{sec:data-limits}: the labels test pose,
gender, height and proportions but not body build, and no number compares with published 3DPW results. A body-joint
MPJPE next to the \MethodNumBones-joint one would help, but is not written.

\paragraph{2a (met).} His expectation held: the real images, not the fixes, brought the large 3DPW gain. Nothing is
missing.

\paragraph{2b (not started).} No temporal code exists. Section~\ref{sec:future} plans a small refiner on the frozen
Final model; whether a refiner on frozen features counts as the temporal model he means is a question for him.

\paragraph{3 (partly).} No checkpoint is released, the scripts still carry the cluster's paths, and the licences have
not been checked (Section~\ref{sec:disc-limitations}). Section~\ref{sec:future} lists the release steps.

% -----------------------------------------------------------------------------
\section{Why PVE barely moved while everything else did}
\label{sec:disc-pve}

From Control to Final the root error fell by \RootErrDropPct\,\%, neck02 from \ControlNeckTwo{} to
\FinalNeckTwo\,deg and the finger error from \ControlFingers{} to \FinalFingers\,deg, while PVE fell only from
\ControlPVE{} to \FinalPVE\,mm, within run-to-run variation, and Base keeps the lowest PVE
(Table~\ref{tab:final-annyone}). Three facts explain most of this.

\paragraph{PVE cannot see placement.} PVE aligns the pelvises before it measures, so root error, depth and drawn size
are removed by construction; on Base the correlation between a person's PVE and their depth error was
\DiagCorrPVEDepth{} (Table~\ref{tab:diag-placement}).

\paragraph{The bone flips cancelled.} The half-turns of neck02 and the head bone cancel, so Base's head pointed the
right way and only the neck skin between them was twisted: \DiscNeckVerts{} of \num{\MethodNumVerts} vertices. Its
error fell from \ControlNeckSkin{} to \FinalNeckSkin\,mm, about \DiscNeckPVEShare\,mm spread over the whole mesh (an
estimate, not a measurement).

\paragraph{Noise floor and restart cost.} One run resolves about \MethodResolvePVE\,mm PVE
(Section~\ref{sec:diag-noise}), so the gap of \GapPVE\,mm supports neither a gain nor a cost. Restarting a fine-tune
of Base already cost about \RestartCost\,mm, and Control sits \DiscBaseControlGapPVE\,mm above Base, likely because
every fine-tune moves away from the synthetic domain of the holdout; this was not tested.

What did move is PA-PVE, which also removes global rotation and scale: Final is \GapPAPVE\,mm below Control and
\DiscPAPVEBelowBase\,mm below Base, and the pilots tie part of this to the bone correction (Table~\ref{tab:neck-bones}).
The likely reason is that the error in global rotation dominates PVE and hides the articulated part where the flips
sit; a per-region breakdown would test this. I therefore report each fix with the measure that can see it, and keep
PVE in every table as the check that a fix costs nothing a single run can resolve.

% -----------------------------------------------------------------------------
\section{Limitations}
\label{sec:disc-limitations}

\paragraph{Children: a measured improvement, not solved.} Only \RHBabiesBelowFinal\,\% of the babies and
\RHKidsBelowFinal\,\% of the kids fall below the child line, and the boy of the demo photos gets an adult body.
Children are rare in AnnyOne and absent from the holdout, and the 3DPW labels pin every real person's age at
\DataFixedAge. Whether the domain gap or an L1 loss through a sigmoid on a rare value dominates, the evidence cannot
rank.

\paragraph{Depth: the slope did not move.} The fix moved the offset and the size, not the slope (\FinalSlope):
\FinalTooFarPct\,\% of all people, the near ones most, are still placed too far. The cause is unknown, and all depth
numbers are synthetic.

\paragraph{One run per arm, no seed.} The claims rest on gaps far above the run-to-run variation (root error, bone
angles, girth, fingers, PA-PVE); the PVE gap, the 3DPW gain of the fixes and the recall cost lie within or at its edge.
There are no bootstrap intervals, and the holdout holds \DataHoldoutImages{} images with \DiagHoldoutPeople{} people.
The finger gain, likewise, is measured on the synthetic holdout only; on real images it rests on three images by eye.

\paragraph{The recall cost.} Final finds \GapRecallBase{} points fewer people than Base, beyond the variation, and every
depth-fix pilot lost recall against its control. I state it as a real cost; the head-anchored detection cell is the
likely reason, not yet tested.

\paragraph{What the 3DPW labels can teach and test.} The labels carry a neutral body, with \DataConstantCount{} of
\DataPhenotypes{} phenotypes constant and the neck, head and hand bones at rest (Section~\ref{sec:data-limits}). They
cannot supervise body build or fingers or test the neck correction, they put a floor of about \DataFloor\,mm under
every number, and precision is a lower bound because bystanders count as false positives.

\paragraph{The comparison with NAVER's model.} Its AnnyOne number is most likely in-sample
(Section~\ref{sec:res-baseline}). On 3DPW, Final was fine-tuned on 3DPW-train and NAVER's model was not; NAVER's model
fine-tuned on our labels would separate network from data and was not run.

\paragraph{Licences.} SMPL and 3DPW are licensed for non-commercial research, Relative Human ships without a licence
file (I used only its validation split, for evaluation), and AnnyOne and NAVER's checkpoints carry their own terms. I
have not yet checked whether these allow releasing derived weights.

% -----------------------------------------------------------------------------
\section{Engineering lessons}
\label{sec:disc-engineering}

The costly faults were silent: each gave plausible output and showed only when one stage was measured on its own.
Four habits found them. I measured each stage alone; only an independent check of the finished label files found the
three label defects (Section~\ref{sec:data-defects}). I read every change against a control, which kept the neck
pilot's PVE and the 3DPW gain from being credited to the fixes. I checked by geometry, not by names: the reversed
spine, the reversed gender value, the pelvis targets and the half-turned neck bones all hid behind plausible names or
group means. And every change entered the code as a switch whose default keeps the old behaviour
(Appendix~\ref{app:flags}).
